\chapter[\emj{⚖️}~AVL Trees (Árboles Binarios Balanceados)]{\emj{⚖️}~AVL Trees (Árboles Binarios Balanceados)}

\begin{dsaquees}

Un móvil de bebé 👶. Si pones mucho peso de un lado, el móvil se
inclina y queda chueco. Para que funcione bien, debes mover las piezas
hasta que quede perfectamente equilibrado.

Un árbol AVL es un Binary Search Tree (BST) que tiene una "manía"
por el orden: no permite que una rama sea mucho más larga que otra. 
Si al insertar un número el árbol se desbalancea, el árbol se 
"tuerce" a sí mismo (rotación) para volver al equilibrio.

\end{dsaquees}

\begin{dsaotraforma}

El AVL es el primer árbol auto-balanceable de la historia. Su regla
es estricta: para cada nodo, la diferencia de altura entre su hijo
izquierdo y derecho (llamada Balance Factor) solo puede ser -1, 0 o 1.

Si la diferencia es 2 o -2, el árbol realiza rotaciones (Simples o 
Dobles) para restaurar el balance. Esto garantiza que la altura del
árbol sea siempre O(log n), evitando el peor caso de los BST normales.

\end{dsaotraforma}

\begin{dsadiagrama}
Rotación Simple a la Derecha (LL Rotation):

      [3] (BF=+2)              [2] (BF=0)
     /                        /   \
   [2]          ----->      [1]    [3]
  /
[1]

Rotación Doble (LR Rotation):

    [3] (BF=+2)         [3]                 [2]
   /                   /                   /   \
 [1] (BF=-1)  --->   [2]      --->       [1]    [3]
   \                /
   [2]            [1]
\end{dsadiagrama}

\begin{lstlisting}
1. Insert: Igual que un BST, pero después de insertar, el árbol revisa
   la altura de los nodos hacia arriba. Si encuentra uno desbalanceado,
   aplica una de las 4 rotaciones (LL, RR, LR, RL).
2. Delete: También igual que BST, pero requiere revisar el balance 
   hacia arriba, pudiendo causar múltiples rotaciones.
3. Search: Exactamente igual que en un BST, pero con la garantía de
   que siempre será rápido (O(log n)).

PSEUDOCÓDIGO
──────────────────────────────────────────────────────────────

// Obtener factor de balance
función obtenerBalance(nodo):
    SI nodo es NULL: return 0
    return altura(nodo.izq) - altura(nodo.der)

// Rotación a la derecha
función rotarDerecha(y):
    x = y.izq
    T2 = x.der
    x.der = y
    y.izq = T2
    actualizarAlturas(y, x)
    return x

CÓDIGO C++
──────────────────────────────────────────────────────────────

#include <iostream>
#include <algorithm>
using namespace std;

struct Node {
    int key, height;
    Node *left, *right;
    Node(int k) : key(k), height(1), left(nullptr), right(nullptr) {}
};

int getHeight(Node* n) { return n ? n->height : 0; }

Node* rightRotate(Node* y) {
    Node* x = y->left;
    Node* T2 = x->right;
    x->right = y;
    y->left = T2;
    y->height = max(getHeight(y->left), getHeight(y->right)) + 1;
    x->height = max(getHeight(x->left), getHeight(x->right)) + 1;
    return x;
}

Node* leftRotate(Node* x) {
    Node* y = x->right;
    Node* T2 = y->left;
    y->left = x;
    x->right = T2;
    x->height = max(getHeight(x->left), getHeight(x->right)) + 1;
    y->height = max(getHeight(y->left), getHeight(y->right)) + 1;
    return y;
}

Node* insert(Node* node, int key) {
    if (!node) return new Node(key);
    if (key < node->key) node->left = insert(node->left, key);
    else if (key > node->key) node->right = insert(node->right, key);
    else return node;

    node->height = 1 + max(getHeight(node->left), getHeight(node->right));
    int balance = getHeight(node->left) - getHeight(node->right);

    if (balance > 1 && key < node->left->key) return rightRotate(node);
    if (balance < -1 && key > node->right->key) return leftRotate(node);
    if (balance > 1 && key > node->left->key) {
        node->left = leftRotate(node->left);
        return rightRotate(node);
    }
    if (balance < -1 && key < node->right->key) {
        node->right = rightRotate(node->right);
        return leftRotate(node);
    }
    return node;
}

int main() {
    Node* root = nullptr;
    root = insert(root, 10);
    root = insert(root, 20);
    root = insert(root, 30); // Esto causará una rotación RR

    cout << "Raíz del árbol balanceado: " << root->key << endl; // Debería ser 20
    return 0;
}

COMPLEJIDAD (Big-O)
──────────────────────────────────────────────────────────────

  Operación   │ Tiempo Garantizado
  ────────────┼───────────────────
  Search      │ O(log n)
  Insert      │ O(log n)
  Delete      │ O(log n)
  Espacio     │ O(n)

* A diferencia del BST normal, el peor caso en AVL sigue siendo O(log n).
\end{lstlisting}

\begin{lstlisting}
✅ ¿Por qué AVL y no BST?: La respuesta siempre es "Balance Garantizado".
   Un AVL nunca se convertirá en una Linked List.

💡 AVL vs Red-Black Tree: 
   - AVL es más balanceado (búsqueda más rápida).
   - Red-Black es más rápido para insertar/eliminar (menos rotaciones).
   - `std::map` en C++ suele usar Red-Black.

⚠️ Implementación: Casi nunca te pedirán programar un AVL completo 
   en una entrevista de 45 min (es muy largo), pero DEBES saber explicar 
   las rotaciones y por qué son necesarias.
\end{lstlisting}

\begin{dsaresumen}

• BST que se auto-balancea solo.
• Balance Factor: diferencia de altura entre hijos (máximo 1).
• Rotaciones: El truco místico para equilibrar el árbol.
• Garantiza O(log n) para todas las operaciones.
• Más complejo de programar que un BST normal.
• Ideal para aplicaciones donde hay MUCHAS búsquedas y pocas inserciones.
════════════════════════════════════════════════════════════════

\end{dsaresumen}


